\section{Small support on one side of a six-terminal sector}
\label{sec:path-six-matching-support}

Two six-terminal matchings are compatible precisely when they have no common pair.
The next finite observation reduces the possible supports of two compatible forest profiles.

\begin{lemma}[Three matchings have at most two common partners]
\label{lem:path-six-three-common-partners}
Let $M_1,M_2,M_3$ be distinct perfect matchings on six vertices.
At most two perfect matchings have no pair in common with any of them.
\end{lemma}

\begin{proof}
Distinct perfect matchings on six vertices have either zero or one common pair.
Let $t$ be the sum of these three pairwise intersection sizes.

If $t=0$, the union of $M_1,M_2,M_3$ is a simple cubic graph.
Its complement in $K_6$ is a simple two-regular graph.
It is either a six-cycle or two triangles.
These graphs have two or zero perfect matchings, respectively.

If $t=3$, the three matchings have a common pair.
Indeed, three distinct pairwise shared pairs would be disjoint and would force all three matchings to be equal.
The remaining four vertices then use all three perfect matchings of $K_4$.
The allowed-edge graph is $K_{2,4}$ and has no perfect matching.

It remains to check $t=1,2$.
Relabel a pair of intersecting matchings as
\[
 M_1=01|23|45,\qquad M_2=01|24|35.
\]
For $t=2$, relabel the one additional shared pair as $45$.
The two possible pairings on the other four vertices are equivalent under the symmetries of $M_1,M_2$.
For $t=1$, the third matching avoids both displayed matchings.
It must use one of $25,34$, with its two other pairs joining $0,1$ to the remaining vertices.
These four choices are also equivalent under those symmetries.
Thus the complete cases are
\[
\begin{array}{c|c|c}
 t&M_3&\text{common compatible matchings}\\\hline
 2&02|13|45&03|14|25,\quad05|12|34\\
 1&02|15|34&03|14|25,\quad04|13|25
\end{array}
\]
The last column follows by choosing the partner of vertex $0$ among its allowed neighbors, then completing the matching.
Each case has exactly two choices.
\end{proof}

\begin{corollary}[One support has size at most two]
\label{cor:path-six-small-cross-support}
Let $\mathcal A,\mathcal B$ be nonempty families of six-terminal perfect matchings.
Suppose every matching in $\mathcal A$ is compatible with every matching in $\mathcal B$.
Then
\[
 \min\{|\mathcal A|,|\mathcal B|\}\le2.
\]
\end{corollary}

\begin{proof}
If $\mathcal A$ contains three distinct matchings, Lemma~\ref{lem:path-six-three-common-partners} gives $|\mathcal B|\le2$.
Otherwise $|\mathcal A|\le2$ already.
\end{proof}

For a critical pair of seven-pole profiles, the positive matching supports on an active six-terminal sector satisfy this corollary.
One shore therefore uses one or two matchings on that sector.
With one fixed matching, the other shore has at most eight possible matching states.
With two fixed distinct matchings, it has at most four.
This is a support reduction, not a bound on the number of active four-terminal sectors.
